% Part: computability
% Chapter: computability-theory
% Section: applications-fixed-point

\documentclass[../../../include/open-logic-section]{subfiles}

\begin{document}

\olfileid{cmp}{thy}{apf}
\olsection{નિશ્ચિત-બિંદુના પ્રમેયનો ઉપયોગ}

નિશ્ચિત-બિંદુનું પ્રમેય આંશિક સંગણનીય વિધેયોને તેમના
સૂચકોના આધારે વ્યાખ્યાયિત કરવાની છૂટ આપે છે. દાખલા
તરીકે, એવો સૂચક $e$ મળે છે કે દરેક $y$ માટે
\[
\cfind{e}(y) = e + y.
\]
બીજા ઉદાહરણ તરીકે, નિશ્ચિત-બિંદુના પ્રમેયની સાબિતીથી
પોતાને જ છાપતો Java અથવા C++ પ્રોગ્રામ ઘડી શકાય છે.

યાદ કરો કે દરેક $e$ માટે $W_e$ને $\cfind{e}$નો પ્રદેશ
રાખીએ, તો અનુક્રમ $W_0$, $W_1$, $W_2$,~\dots સંગણકીય
રીતે પરિગણનીય ગણોની પરિગણના કરે છે. આ ગણોમાંના કેટલાક
સંગણનીય છે. એવો પ્રશ્ન થાય કે શું એવી કલનવિધિ છે જે
મૂલ્ય $x$ ઇનપુટ તરીકે લે અને, જો $W_x$ સંગણનીય હોય,
તો તેના લક્ષણ વિધેયનો સૂચક પાછો આપે? જવાબ ``ના'' છે;
આવી કલનવિધિ નથી:

\begin{thm}
આગળનો ગુણધર્મ ધરાવતું કોઈ આંશિક સંગણનીય વિધેય $f$
નથી: જ્યારે $W_e$ સંગણનીય હોય, ત્યારે $f(e)$
વ્યાખ્યાયિત હોય અને $\cfind{f(e)}$ તેનું લક્ષણ વિધેય હોય.
\end{thm}

\begin{proof}
$f$ કોઈ પણ આંશિક સંગણનીય વિધેય હોય; આપણે એવો $e$
બનાવીશું કે $W_e$ સંગણનીય હોય, પરંતુ $\cfind{f(e)}$
તેનું લક્ષણ વિધેય ન હોય. નિશ્ચિત-બિંદુના પ્રમેયથી
એવો સૂચક $e$ મળે છે કે
\[
\cfind{e}(y) \simeq
\begin{cases}
0 & \text{જો $y=0$ અને $\cfind{f(e)}(0) \fdefined = 0$} \\
\text{અનિર્ધારિત} & \text{અન્યથા.}
\end{cases}
\]
એટલે $e$ નીચેના વિધેય પર નિશ્ચિત-બિંદુનું પ્રમેય
લાગુ કરીને મળે છે:
\[
g(x,y) \simeq
\begin{cases}
0 & \text{જો $y=0$ અને $\cfind{f(x)}(0) \fdefined = 0$} \\
\text{અનિર્ધારિત} & \text{અન્યથા.}
\end{cases}
\]
અનૌપચારિક રીતે $g$ આંશિક સંગણનીય છે: ઇનપુટ $x$
અને $y$ પર કલનવિધિ પહેલાં તપાસે છે કે $y$ એ~$0$
છે કે નહીં. હોય, તો તે $f(x)$ ગણે છે અને પછી
સર્વવ્યાપી મશીનથી $\cfind{f(x)}(0)$ ગણે છે. આ
છેલ્લી સંગણના થંભીને~$0$ આપે, તો કલનવિધિ~$0$ આપે;
અન્યથા તે થંભતી નથી.
\sourcecorrection{OLCT-012}{સ્થિર અંગ્રેજી મૂળ સાબિતી
શરૂ કરતાં $f$ને સંપૂર્ણ સંગણનીય માને છે, જ્યારે પ્રમેય
બધાં આંશિક સંગણનીય વિધેયોને નકારે છે. ઉપરના $g$ની
ગણતરી $f(x)$ અનિર્ધારિત રહે ત્યારે પણ આંશિક સંગણનીય
રહે છે; તેથી સાબિતી આ વ્યાપક દાવાને પણ આવરે છે.}

હવે નોંધો કે $\cfind{f(e)}(0)$ વ્યાખ્યાયિત થઈને $0$
જેટલું હોય, તો $\cfind{e}(y)$ ત્યારે અને માત્ર ત્યારે
વ્યાખ્યાયિત છે જ્યારે $y$, $0$ હોય; તેથી $W_e =
\{ 0 \}$. જો $\cfind{f(e)}(0)$ અનિર્ધારિત હોય અથવા
વ્યાખ્યાયિત હોય પરંતુ $0$ ન હોય, તો $W_e = \emptyset$.
બંને રીતે $\cfind{f(e)}$, $W_e$નું લક્ષણ વિધેય નથી:
ઇનપુટ $0$ પર તે ખોટો જવાબ આપે છે. જો સૂચક આપતું
વિધેય જ આ ઇનપુટ પર અનિર્ધારિત હોય, તો માગેલી શરત
પહેલેથી જ ભંગાય છે.
\end{proof}

\end{document}
